\input{../../../templates/es/preambulo.tex}

% Los listados se toman de src/ con \lstinputlisting; la diapositiva es el archivo.
\lstset{
    basicstyle=\ttfamily\fontsize{8pt}{9.6pt}\selectfont,
    breaklines=true,
    breakatwhitespace=true,
    columns=fullflexible,
    keepspaces=true,
    xleftmargin=0.3em,
    framesep=2pt,
    aboveskip=0pt,
    belowskip=0pt
}

\newcommand{\marcoresultado}[3]{%
    \begin{frame}{#1}
        \begin{center}
            \includegraphics[height=0.62\textheight,width=\textwidth,keepaspectratio]{#2}
        \end{center}
        \vspace{0.25em}
        {\small #3\par}
    \end{frame}%
}

\newcommand{\marcoafirmacion}[2]{%
    \begin{frame}{#1}
        {\small #2\par}
    \end{frame}%
}

\date{}

\begin{document}

\tituloleccion{02}{El flujo del algoritmo genético}

\section{Esta lección}

\begin{frame}{Esta lección}
El paso 6 de la Lección 01, con el ciclo todavía sin nombre, se reproduce aquí al dígito. Nombrar el flujo no debe cambiar la respuesta.

\vspace{0.6em}
\begin{itemize}
    \item Cinco fases con nombre; una generación es una función.
    \item El problema es un argumento, no va soldado al conductor.
    \item Una regla de parada barata reporta que la ejecución dejó de mejorar --- no que terminó.
\end{itemize}
\end{frame}

% ================= flujo_ag_01_las_cinco_fases =================
\begin{frame}{Modelo matemático}
\[
P_{t+1}=R\!\left(M\!\left(C\!\left(S(P_t)\right)\right)\right)
\]
Una generación es una transición aleatoria de poblaciones. Una prueba de
paciencia como \(b_t-b_{t-k}<\varepsilon\) detecta poca mejora reciente; no
prueba convergencia ni dice nada del óptimo global desconocido.

Para \(r\lfloor(T-s|x-x_0|)/h\rfloor\), las bandas unilaterales ordinarias
miden \(h/s=0.5\); aquí la cima es \(|x-4|\le0.25\), ancho 0.5 y no un punto.
\end{frame}

\section{El flujo, nombrado}

\begin{frame}{El flujo, nombrado}
La Lección 01 terminó con un algoritmo genético que funcionaba. Su cuerpo de ciclo, sin embargo,
era una serie de bloques de comentarios sin nombre, y no puedes discutir "el flujo" cuando el
flujo no tiene nombre y no tiene límite. Este paso lo nombra.
\end{frame}

\begin{frame}[fragile]{(1) evolve\_one\_generation()}
\lstinputlisting[basicstyle=\ttfamily\fontsize{6pt}{7.2pt}\selectfont,firstline=94,lastline=127]{../src/flujo_ag_01_las_cinco_fases.py}
\end{frame}

\begin{frame}[fragile]{(2) run()}
\lstinputlisting[basicstyle=\ttfamily\fontsize{7pt}{8.5pt}\selectfont,firstline=130,lastline=148]{../src/flujo_ag_01_las_cinco_fases.py}
\end{frame}

\begin{frame}[fragile]{(3) el rastro de fases}
\lstinputlisting[basicstyle=\ttfamily\fontsize{6pt}{7.2pt}\selectfont,firstline=153,lastline=184]{../src/flujo_ag_01_las_cinco_fases.py}
\end{frame}

\marcoresultado{El flujo, nombrado}{../figures/flujo_ag_01_las_cinco_fases.png}{%
    Misma semilla, misma respuesta que la Lección 01: \texttt{x=+1.372 f=+0.706}. La ejecución costó \textbf{107} evaluaciones de aptitud, no las 110 obvias — un individuo seleccionado pero no cruzado ni mutado es el mismo objeto y nunca se vuelve a evaluar}

% ================= flujo_ag_02_aptitud_inyectada =================
\section{El problema se convierte en un argumento}

\begin{frame}{El problema se convierte en un argumento}
En el paso 1, la palabra \texttt{objective} aparece dentro de Individuo. Esa única referencia
suelda el flujo a un problema: para resolver un segundo problema copiarías el archivo.
El libro separa \texttt{fitness.py} de \texttt{individual.py} exactamente por esta razón.
\end{frame}

\begin{frame}[fragile]{(1) los dos problemas}
\lstinputlisting[basicstyle=\ttfamily\fontsize{8pt}{9.6pt}\selectfont,firstline=47,lastline=61]{../src/flujo_ag_02_aptitud_inyectada.py}
\end{frame}

\begin{frame}[fragile]{(2) Individuo(genes, f)}
\lstinputlisting[basicstyle=\ttfamily\fontsize{8pt}{9.6pt}\selectfont,firstline=69,lastline=82]{../src/flujo_ag_02_aptitud_inyectada.py}
\end{frame}

\begin{frame}[fragile]{(3) los operadores de construcción}
\lstinputlisting[basicstyle=\ttfamily\fontsize{8pt}{9.6pt}\selectfont,firstline=99,lastline=115]{../src/flujo_ag_02_aptitud_inyectada.py}
\end{frame}

\begin{frame}[fragile]{(4) run(fitness\_function)}
\lstinputlisting[basicstyle=\ttfamily\fontsize{7pt}{8.5pt}\selectfont,firstline=147,lastline=165]{../src/flujo_ag_02_aptitud_inyectada.py}
\end{frame}

\begin{frame}[fragile]{(5) las dos ejecuciones}
\lstinputlisting[basicstyle=\ttfamily\fontsize{7pt}{8.5pt}\selectfont,firstline=168,lastline=190]{../src/flujo_ag_02_aptitud_inyectada.py}
\end{frame}

\marcoresultado{El problema se convierte en un argumento}{../figures/flujo_ag_02_aptitud_inyectada.png}{%
    Un controlador, dos problemas: la ejecución del seno reproduce exactamente el paso 1, la ejecución del objetivo aterriza a \textbf{0.0038} de +4.200. La única diferencia entre las ejecuciones es qué función se pasó}

% ================= flujo_ag_03_metricas_poblacion =================
\section{Midiendo la población, no al líder}

\begin{frame}{Midiendo la población, no al líder}
Hasta ahora, cada rastro reporta el mejor individuo. Ese es el único número sobre
el que un AG es menos honesto: el mejor puede quedarse quieto durante cinco
generaciones mientras la población debajo de él cambia por completo, o puede ser un
único caso atípico afortunado en una población que no ha aprendido nada.
\end{frame}

\begin{frame}[fragile]{(1) population\_metrics()}
\lstinputlisting[basicstyle=\ttfamily\fontsize{8pt}{9.6pt}\selectfont,firstline=67,lastline=83]{../src/flujo_ag_03_metricas_poblacion.py}
\end{frame}

\begin{frame}[fragile]{(2) la tabla de historial}
\lstinputlisting[basicstyle=\ttfamily\fontsize{6pt}{7.2pt}\selectfont,firstline=188,lastline=230]{../src/flujo_ag_03_metricas_poblacion.py}
\end{frame}

\begin{frame}[fragile]{(3) la figura de métricas}
\lstinputlisting[basicstyle=\ttfamily\fontsize{7pt}{8.5pt}\selectfont,firstline=232,lastline=257]{../src/flujo_ag_03_metricas_poblacion.py}
\end{frame}

\marcoresultado{Midiendo la población, no al líder}{../figures/flujo_ag_03_metricas_poblacion.png}{%
    En el problema del seno, la aptitud promedio sube de −1.5588 a +0.6228 mientras que la dispersión genética cae de \textbf{6.665 a 0.064} para la generación 8. El campeón se queda quieto mientras la población colapsa debajo de él}

% ================= flujo_ag_04_condicion_parada =================
\section{Sabiendo cuándo detenerse}

\begin{frame}{Sabiendo cuándo detenerse}
Cada ejecución hasta ahora se detuvo después de exactamente diez generaciones, porque
diez estaba escrito en la parte superior del archivo. El libro enumera tres formas de
terminar una ejecución: se encuentra una solución aceptable, se alcanza un conteo de
generaciones, o las mejoras se agotan. El paso 3 midió exactamente cuándo se agotan
las mejoras, por lo que la tercera condición ahora está disponible gratis.
\end{frame}

\begin{frame}[fragile]{(1) MAX\_GENERATIONS}
\lstinputlisting[basicstyle=\ttfamily\fontsize{8pt}{9.6pt}\selectfont,firstline=40,lastline=40]{../src/flujo_ag_04_condicion_parada.py}
\end{frame}

\begin{frame}[fragile]{(2) PATIENCE, MIN\_IMPROVEMENT}
\lstinputlisting[basicstyle=\ttfamily\fontsize{8pt}{9.6pt}\selectfont,firstline=45,lastline=52]{../src/flujo_ag_04_condicion_parada.py}
\end{frame}

\begin{frame}[fragile]{(3) detenerse por estancamiento}
\lstinputlisting[basicstyle=\ttfamily\fontsize{8pt}{9.6pt}\selectfont,firstline=186,lastline=195]{../src/flujo_ag_04_condicion_parada.py}
\end{frame}

\begin{frame}[fragile]{(4) las ejecuciones de control}
\lstinputlisting[basicstyle=\ttfamily\fontsize{6pt}{7.2pt}\selectfont,firstline=274,lastline=327]{../src/flujo_ag_04_condicion_parada.py}
\end{frame}

\marcoresultado{Sabiendo cuándo detenerse}{../figures/flujo_ag_04_condicion_parada.png}{%
    La regla ahorra de \textbf{49 a 51} de las 60 generaciones permitidas y cuesta menos que \texttt{MIN\_IMPROVEMENT} en ambos problemas. También expone un defecto: en 1 de 2 problemas, la última generación ya no contiene al campeón — nada aquí lo protege}

% ================= flujo_ag_05_paisajes_escalonados =================
\section{Paisajes escalonados, y una suposición que resultó errónea}

\begin{frame}{Paisajes escalonados, y una suposición que resultó errónea}
El paso 4 evaluó el costo de la regla de parada en dos paisajes suaves y encontró que era gratis. La
sospecha natural es que la suavidad estaba haciendo el trabajo: en un paisaje donde el
progreso llega en saltos en lugar de pendientes, un tramo plano significaría que la
población aún no ha encontrado el siguiente escalón, no que la búsqueda haya terminado - y una
regla de paciencia no puede distinguirlos.
\end{frame}

\begin{frame}[fragile]{(1) escalera()}
\lstinputlisting[basicstyle=\ttfamily\fontsize{6pt}{7.2pt}\selectfont,firstline=81,lastline=112]{../src/flujo_ag_05_paisajes_escalonados.py}
\end{frame}

\begin{frame}[fragile]{(2) PROBLEMAS}
\lstinputlisting[basicstyle=\ttfamily\fontsize{8pt}{9.6pt}\selectfont,firstline=238,lastline=242]{../src/flujo_ag_05_paisajes_escalonados.py}
\end{frame}

\begin{frame}[fragile]{(3) el veredicto}
\lstinputlisting[basicstyle=\ttfamily\fontsize{6pt}{7.2pt}\selectfont,firstline=368,lastline=439]{../src/flujo_ag_05_paisajes_escalonados.py}
\end{frame}

\marcoresultado{Paisajes escalonados, y una suposición que resultó errónea}{../figures/flujo_ag_05_paisajes_escalonados.png}{%
    La escalera \textbf{no} rompe la regla: la ejecución llega a \texttt{+2.0000}, la referencia de cuadrícula densa del paisaje, y la regla la detiene sin haber perdido nada. Cada uno de los tres problemas coincide con su referencia de cuadrícula, déficit \texttt{+0.0000}}

\section{Siguiente}

\begin{frame}{Siguiente}
En la escalera la regla de paciencia nunca se rinde demasiado pronto: los peldaños suben, y cada borde es una comparación que la selección puede usar. La hipótesis queda refutada, y se conserva.

\vspace{0.8em}
\textbf{Lección 03 --- Selección}

El primer operador que se abre, todavía sobre una generación de una población. La evidencia a nivel de corrida se aplaza a la Lección 06.
\end{frame}
\end{document}
